\chapter{Conclusion and Outlook}

\section{Conclusion}

This work presented a measurement-based comparison of Ascon-AEAD128 and
AES-128-GCM on a Cortex-M0 and a Cortex-M4F, using software-only implementations
under matched conditions.
It was motivated by the absence of quantitative performance data for the
standardized lightweight algorithm on representative constrained hardware, and it
set out to provide that data across execution time, code size, optimization
level, and processor architecture.

The measurements answer the four research questions consistently: Ascon-AEAD128
was faster at every message size and optimization level on both boards, required
about one quarter of the Flash of AES-128-GCM at the size-optimized level, showed
an advantage that was understated at \texttt{-O0} and clearer once optimization
was enabled, and showed an advantage that grew on the more constrained
Cortex-M0.
Taken together the results give a quantitative basis for preferring
Ascon-AEAD128 on constrained, software-only targets, subject to the
qualification that the speed advantage and the code-size advantage are obtained
together only when building for size.

\section{Outlook}

Several directions extend this work.
First, measuring the larger-table AES-128-GCM configuration would bound the
fastest AES-128-GCM throughput achievable on these boards and quantify the speed
gained for the additional RAM.
Second, running the Cortex-M4F at its native 80\,MHz would give wall-clock
figures for that board at full speed, confirming the cycle-count results that
are independent of operating frequency.
Third, comparing Ascon-Hash256 against SHA-256 on the same hardware would extend
the comparison to the hashing use case that arises in key derivation and
integrity checking.
Fourth, repeating the study on a part with a hardware AES accelerator would
characterize how that shifts the balance and would mark the boundary of the
recommendation made here.
Fifth, measuring Ascon-80pq, which shares the same permutation as
Ascon-AEAD128~\cite{ASCONSpec}, would quantify the cost of its longer key and
position the comparison against longer-term security requirements.
These directions all build on the measurement infrastructure established here,
which was validated, frozen, and documented for reproducibility.
